%! Unidad 13 — Termodinámica: ejercicios resueltos, ejercicios finales y autoevaluación
%! Fuente: Fuentes/Fisica/Unidad 13 («Ejercicios resueltos», «Practica y calcula», «Piensa y resuelve»). La fuente no trae
%! olimpiadas ni autoevaluación de este tema: la autoevaluación es propia. Enunciados adaptados con datos cambiados y soluciones
%! rehechas con 1 cal = 4,18 J, R = 8,31 J/(mol·K) y 1 atm = 101 325 Pa.
%! Dificultad: \(\star\) baja, \(\star\star\) media, \(\star\star\star\) alta.

\section{Ejercicios resueltos}\label{sec:ud13-resueltos}

\begin{example}
    \textbf{Los cilindros de una moto.} El motor de una moto de competición tiene 4 cilindros de \qty{250}{\cm\cubed} cada uno. En
    cada ciclo, la mezcla de aire y combustible se comprime hasta ocupar \qty{25}{\cm\cubed}, y se pierden \qty{15}{\joule} de calor
    por cilindro. Suponiendo (para simplificar) que la compresión se produce a una presión constante de \qty{101}{\kilo\pascal},
    calcula: a) el trabajo que hace cada cilindro sobre el gas; b) la variación de energía interna del gas.

    \textbf{Solución.} \textbf{1. Unidades SI}: \(V_i = \qty{2.5e-4}{\m\cubed}\), \(V_f = \qty{2.5e-5}{\m\cubed}\),
    \(p = \qty{1.01e5}{\pascal}\).

    \textbf{2. a)} Es una compresión:
    \(W = -p\,(V_f - V_i) = -\num{1.01e5} \cdot (\num{2.5e-5} - \num{2.5e-4}) = \qty{+22.7}{\joule}\). Es positivo porque el
    entorno (el pistón) hace trabajo sobre el gas.

    \textbf{3. b)} El gas cede calor, \(Q = \qty{-15}{\joule}\):
    \(\Delta U = Q + W = -15 + \num{22.7} = \qty{+7.7}{\joule}\). La energía interna (y la temperatura) de la mezcla aumenta.
\end{example}

\begin{example}
    \textbf{Mezcla de agua con paredes adiabáticas y diatérmanas.} En un recipiente con \qty{1500}{\cm\cubed} de agua a
    \qty{12}{\degreeCelsius} se echan \qty{500}{\cm\cubed} de agua a \qty{72}{\degreeCelsius}. La densidad del agua es
    \qty{1}{\g\per\cm\cubed} y \(c_e = \qty{1}{\calorie\per\g\per\degreeCelsius}\). a) Si las paredes son adiabáticas, ¿cuál es la
    temperatura final? b) ¿Qué calor intercambia cada parte? c) ¿Qué ocurriría si las paredes fueran diatérmanas y el exterior
    estuviera a \qty{22}{\degreeCelsius}?

    \textbf{Solución.} \textbf{1. Masas}: \(m_1 = \qty{1500}{\g}\) (fría) y \(m_2 = \qty{500}{\g}\) (caliente).

    \textbf{2. a)} Con paredes adiabáticas, \(Q_1 + Q_2 = 0\):
    \[ m_1\,c_e\,(T - T_1) + m_2\,c_e\,(T - T_2) = 0 \Rightarrow T = \frac{m_1\,T_1 + m_2\,T_2}{m_1 + m_2} = \frac{1500 \cdot 12 +
    500 \cdot 72}{2000} = \qty{27}{\degreeCelsius} \]
    No hace falta pasar a kelvin, porque solo intervienen diferencias de temperatura.

    \textbf{3. b)} \(Q_1 = 1500 \cdot 1 \cdot (27 - 12) = \qty{22500}{\calorie} = \qty{94.1}{\kilo\joule}\) (absorbe);
    \(Q_2 = \qty{-94.1}{\kilo\joule}\) (cede lo mismo).

    \textbf{4. c)} Con paredes diatérmanas, todo el sistema acaba en equilibrio con el exterior, a \qty{22}{\degreeCelsius} (el entorno
    es tan grande que su temperatura no cambia):
    \(Q_1 = 1500 \cdot \num{4.18} \cdot (22 - 12) = \qty{+62.7}{\kilo\joule}\);
    \(Q_2 = 500 \cdot \num{4.18} \cdot (22 - 72) = \qty{-104.5}{\kilo\joule}\). Ahora no se compensan: el sistema cede
    \qty{41.8}{\kilo\joule} al exterior.
\end{example}

\begin{example}
    \textbf{Un ciclo con un tramo isotérmico.} Una masa fija de gas ideal recorre el ciclo de la figura~\ref{fig:ud13-res-ciclo}: de 1
    (\qty{3}{\litre}, \qty{3}{\atm}) a 2 (\qty{6}{\litre}, \qty{1.5}{\atm}) de forma isotérmica; de 2 a 3 (\qty{3}{\litre},
    \qty{1.5}{\atm}) a presión constante, y de 3 a 1 a volumen constante. a) Calcula el trabajo en cada tramo y en el ciclo. b) Calcula
    \(\Delta U\) y \(Q\) en el tramo 1 \(\rightarrow\) 2.

    \textbf{Solución.} \textbf{1. Tramo 1 \(\rightarrow\) 2} (isotérmico). No conocemos \(n\) ni \(T\), pero \(n\,R\,T = p_1\,V_1\):
    \[ W_{1\to2} = -p_1\,V_1 \ln\frac{V_2}{V_1} = -3 \cdot 101325 \cdot \num{3e-3} \cdot \ln 2 = \qty{-632}{\joule} \]

    \textbf{2. Tramo 2 \(\rightarrow\) 3} (isobárico, compresión):
    \(W_{2\to3} = -\num{1.5} \cdot 101325 \cdot (3 - 6) \cdot \num{e-3} = \qty{+456}{\joule}\).

    \textbf{3. Tramo 3 \(\rightarrow\) 1} (isocórico): \(W_{3\to1} = 0\).

    \textbf{4. Ciclo}: \(W = -632 + 456 + 0 = \qty{-176}{\joule}\). El gas realiza trabajo neto sobre el entorno.

    \textbf{5. b)} En el tramo isotérmico de un gas ideal, \(\Delta U_{1\to2} = 0\), y por el primer principio
    \(Q_{1\to2} = -W_{1\to2} = \qty{+632}{\joule}\): el gas absorbe \qty{632}{\joule} de calor y los devuelve como trabajo.
\end{example}

\begin{figure}[H]
    \centering
    \begin{tikzpicture}
        \begin{axis}[grafica, width=0.5\linewidth, height=5cm, xlabel={\(V\) (L)}, ylabel={\(p\) (atm)}, xmin=0, xmax=7.5, ymin=0,
            ymax=3.6, xtick={3,6}, ytick={1.5,3}]
            \addplot[fill=ocre!20, draw=none, domain=3:6] {9/x} \closedcycle;
            \addplot[fill=white, draw=none] coordinates {(3,0) (3,1.5) (6,1.5) (6,0)} \closedcycle;
            \addplot[very thick, ocre, domain=3:6] {9/x};
            \addplot[very thick, ocre, -latex] coordinates {(6,1.5) (4.4,1.5)};
            \addplot[very thick, ocre] coordinates {(4.5,1.5) (3,1.5)};
            \addplot[very thick, ocre, -latex] coordinates {(3,1.5) (3,2.35)};
            \addplot[very thick, ocre] coordinates {(3,2.3) (3,3)};
            \node[font=\small, anchor=south west] at (axis cs:3,3) {1};
            \node[font=\small, anchor=south west] at (axis cs:6,1.5) {2};
            \node[font=\small, anchor=north east] at (axis cs:3,1.5) {3};
            \node[font=\scriptsize, anchor=west] at (axis cs:4.3,2.4) {\(T\) = cte};
        \end{axis}
    \end{tikzpicture}
    \caption{Ciclo con un tramo isotérmico (1 \(\rightarrow\) 2), uno isobárico y uno isocórico.}
    \label{fig:ud13-res-ciclo}
\end{figure}

\begin{example}
    \textbf{Cubitos para el refresco.} Tenemos \qty{330}{\g} de refresco (que trataremos como agua) a \qty{25}{\degreeCelsius}. ¿Qué
    masa de hielo a \qty{0}{\degreeCelsius} hay que añadir para que la mezcla quede a \qty{5}{\degreeCelsius}? Supón que no hay
    intercambio de calor con el exterior. Datos: \(L_f = \qty{80}{\calorie\per\g}\); \(c_{\text{agua}} =
    \qty{1}{\calorie\per\g\per\degreeCelsius}\).

    \textbf{Solución.} \textbf{1. Calor cedido por el refresco}:
    \(Q_1 = 330 \cdot 1 \cdot (5 - 25) = \qty{-6600}{\calorie}\).

    \textbf{2. Calor absorbido por el hielo}: primero se funde a \qty{0}{\degreeCelsius} y después el agua resultante se calienta hasta
    \qty{5}{\degreeCelsius}: \(Q_2 = m\,L_f + m\,c\,(5 - 0) = m\,(80 + 5) = 85\,m\).

    \textbf{3. Balance}: \(Q_1 + Q_2 = 0 \Rightarrow 85\,m = 6600 \Rightarrow m = \qty{77.6}{\g}\), unos cuatro cubitos. Fíjate
    en que la mayor parte del enfriamiento se debe a la fusión del hielo, no a su calentamiento posterior.
\end{example}

\begin{example}
    \textbf{Una central térmica.} En una central térmica, el vapor entra en la turbina a \qty{550}{\degreeCelsius} y sale al
    condensador a \qty{40}{\degreeCelsius}. La central produce \qty{500}{\mega\watt} de potencia eléctrica con un rendimiento real del
    \qty{38}{\percent}. a) ¿Cuál es el rendimiento máximo (de Carnot) entre esas temperaturas? b) ¿Qué potencia térmica toma del
    combustible? c) ¿Qué potencia cede en forma de calor al río o a las torres de refrigeración?

    \textbf{Solución.} \textbf{a)} Temperaturas en kelvin: \(T_1 = \qty{823}{\kelvin}\), \(T_2 = \qty{313}{\kelvin}\).
    \(\eta_{\text{máx}} = 1 - \dfrac{313}{823} = \num{0.62}\), un \qty{62}{\percent}.

    \textbf{b)} \(\eta = \dfrac{P_{\text{útil}}}{P_{\text{absorbida}}} \Rightarrow P_{\text{absorbida}} = \dfrac{500}{\num{0.38}} =
    \qty{1316}{\mega\watt}\).

    \textbf{c)} Por el primer principio (en cada ciclo, \(\Delta U = 0\)), lo que no se convierte en trabajo se cede al foco frío:
    \(P_{\text{cedida}} = 1316 - 500 = \qty{816}{\mega\watt}\). ¡Más de lo que produce! Por eso estas centrales necesitan grandes
    torres de refrigeración o estar junto a un río o al mar.
\end{example}

%========================================================================================
\section{Ejercicios de aplicación}\label{sec:ud13-aplicacion}
%========================================================================================

\subsection{Introducción a la termodinámica}

\begin{exercise}
    \(\star\star\) Explica con tus palabras qué es un sistema termodinámico. Enumera los tipos de sistemas que hay y pon un ejemplo
    de cada uno.
    \begin{solu}
        Es la \textcolor{red}{porción del universo que estudiamos}, separada del entorno por una pared real o imaginaria.
        \textcolor{red}{Abierto} (intercambia materia y energía: una cazuela destapada), \textcolor{red}{cerrado} (solo energía: una
        lata cerrada) y \textcolor{red}{aislado} (ni materia ni energía: un termo ideal).
    \end{solu}
\end{exercise}

\begin{exercise}
    \(\star\) Clasifica en sistemas abiertos, cerrados o aislados: a) el motor de una moto en marcha; b) un móvil cargándose; c) una
    bombona de oxígeno cerrada; d) una célula; e) una lata de sardinas; f) un perro; g) un cactus; h) una ciudad; i) una pelota de
    fútbol; j) un cuenco de sopa caliente.
    \begin{solu}
        \textcolor{red}{Abiertos}: a, d, f, g, h, j. \textcolor{red}{Cerrados}: b, c, e, i. Ninguno es realmente aislado.
    \end{solu}
\end{exercise}

\begin{exercise}
    \(\star\star\) Explica qué es una función de estado. ¿Cuáles de estas magnitudes lo son? a) El trabajo. b) La energía interna. c)
    La entropía. d) El calor. e) La presión. f) El volumen.
    \begin{solu}
        Es una magnitud cuya variación solo depende de los estados inicial y final, no del camino. Son funciones de estado
        \textcolor{red}{b, c, e y f}; el trabajo y el calor \textcolor{red}{no}.
    \end{solu}
\end{exercise}

\begin{exercise}
    \(\star\star\) Razona si son verdaderas o falsas: a) El calor absorbido por un sistema siempre es negativo. b) En un proceso
    isotérmico la temperatura no varía. c) Un proceso a presión constante se llama isocórico. d) En un proceso adiabático hay
    transferencia de calor. e) En un proceso isobárico no cambia el volumen. f) El trabajo que el sistema realiza sobre el entorno
    es positivo.
    \begin{solu}
        a) \textcolor{red}{F}: es positivo. b) \textcolor{red}{V}. c) \textcolor{red}{F}: isobárico. d) \textcolor{red}{F}: \(Q =
        0\). e) \textcolor{red}{F}: lo que no cambia es la presión. f) \textcolor{red}{F}: es negativo (criterio IUPAC).
    \end{solu}
\end{exercise}

\begin{exercise}
    \(\star\) ¿Qué es el equilibrio termodinámico? Si a un sistema en equilibrio le aportamos calor o trabajo, ¿qué le ocurre?
    \begin{solu}
        Es el estado en que las variables del sistema \textcolor{red}{no cambian con el tiempo}. Si le aportamos calor o trabajo,
        \textcolor{red}{se rompe el equilibrio} y el sistema evoluciona hasta alcanzar un nuevo estado de equilibrio.
    \end{solu}
\end{exercise}

\begin{exercise}
    \(\star\) Explica la diferencia entre paredes rígidas y paredes móviles, y entre paredes diatérmanas y adiabáticas.
    \begin{solu}
        Las \textcolor{red}{rígidas} mantienen el volumen constante (no hay trabajo de expansión); las \textcolor{red}{móviles}
        permiten cambiarlo e intercambiar trabajo. Las \textcolor{red}{diatérmanas} dejan pasar el calor y las
        \textcolor{red}{adiabáticas}, no.
    \end{solu}
\end{exercise}

\subsection{El principio cero de la termodinámica}

\begin{exercise}
    \(\star\) ¿Es lo mismo calor que temperatura? Justifica tu respuesta.
    \begin{solu}
        \textcolor{red}{No}. La temperatura es una \textcolor{red}{propiedad del cuerpo} (mide la energía cinética media de sus
        partículas); el calor es \textcolor{red}{energía en tránsito} que pasa de un cuerpo a otro por su diferencia de temperatura.
    \end{solu}
\end{exercise}

\begin{exercise}
    \(\star\) ¿Puede hablarse de equilibrio térmico con un solo cuerpo? Justifica tu respuesta.
    \begin{solu}
        El equilibrio térmico es una relación \textcolor{red}{entre dos o más sistemas} (o entre las partes de un mismo cuerpo):
        se alcanza cuando todos están a la misma temperatura y deja de haber intercambio neto de calor.
    \end{solu}
\end{exercise}

\begin{exercise}
    \(\star\star\) Tres cubos del mismo material y la misma masa, aislados del exterior, se colocan en fila en el orden 3, 1, 2 (el
    1 en el centro, tocando a los otros dos). Sus temperaturas cumplen \(T_3 > T_2 > T_1\). a) ¿Qué cubos ceden calor y cuál lo
    absorbe al principio? b) ¿Cómo acaba el sistema?
    \begin{solu}
        a) El cubo 1, el más frío, \textcolor{red}{absorbe calor de los cubos 3 y 2}, que lo ceden. b) Todos alcanzan la
        \textcolor{red}{misma temperatura}, \(T_{eq} = (T_1 + T_2 + T_3)/3\) (igual masa y material), intermedia entre \(T_1\) y \(T_3\).
    \end{solu}
\end{exercise}

\begin{exercise}
    \(\star\star\) ¿Hay equilibrio térmico? Si no lo hay, ¿cómo evolucionará el sistema para alcanzarlo? a) Al añadir leche fría a
    un café caliente. b) Un termómetro en la axila tras cinco minutos. c) Una botella de agua «del tiempo». d) Un paquete de
    mantequilla recién sacado del frigorífico.
    \begin{solu}
        a) \textcolor{red}{No}: el café cede calor a la leche hasta que la mezcla tiene una temperatura intermedia. b)
        \textcolor{red}{Sí}, con el cuerpo. c) \textcolor{red}{Sí}, con el ambiente. d) \textcolor{red}{No}: absorbe calor del
        ambiente hasta alcanzar su temperatura (y se ablanda).
    \end{solu}
\end{exercise}

\begin{exercise}
    \(\star\) Expresa en kelvin: a) \qty{25}{\degreeCelsius}; b) \qty{-40}{\degreeCelsius}. Expresa en grados Celsius: c)
    \qty{4}{\kelvin} (helio líquido); d) \qty{1500}{\kelvin}.
    \begin{solu}
        a) \(\textcolor{red}{\qty{298.15}{\kelvin}}\). b) \(\textcolor{red}{\qty{233.15}{\kelvin}}\). c)
        \(\textcolor{red}{\qty{-269.15}{\degreeCelsius}}\). d) \(\textcolor{red}{\qty{1226.85}{\degreeCelsius}}\).
    \end{solu}
\end{exercise}

\subsection{Mecanismos de propagación del calor}

\begin{exercise}
    \(\star\) Explica con tus palabras los tres mecanismos de propagación del calor y pon un ejemplo de cada uno.
    \begin{solu}
        \textcolor{red}{Conducción}: por choques entre partículas, sin transporte de materia (una cuchara en la sopa).
        \textcolor{red}{Convección}: por corrientes de fluido caliente que sube y frío que baja (agua al calentarse en una olla).
        \textcolor{red}{Radiación}: por ondas electromagnéticas, incluso en el vacío (el calor del Sol).
    \end{solu}
\end{exercise}

\begin{exercise}
    \(\star\) ¿Por qué nos quemamos si ponemos la mano encima de la llama de una vela, y no si la ponemos al lado a la misma
    distancia?
    \begin{solu}
        Encima nos llega el aire caliente que sube por \textcolor{red}{convección}, además de la radiación; al lado solo nos llega la
        \textcolor{red}{radiación}, que transporta mucha menos energía.
    \end{solu}
\end{exercise}

\begin{exercise}
    \(\star\) ¿Por qué los utensilios de cocina metálicos tienen el mango de plástico o de madera?
    \begin{solu}
        Porque el plástico y la madera son \textcolor{red}{malos conductores del calor} (aislantes): el mango no se calienta y no nos
        quemamos.
    \end{solu}
\end{exercise}

\begin{exercise}
    \(\star\star\) Un termo tiene doble pared de vidrio con vacío entre ellas y las superficies plateadas. Explica qué mecanismo de
    propagación del calor evita cada elemento.
    \begin{solu}
        El \textcolor{red}{vacío} impide la \textcolor{red}{conducción} y la \textcolor{red}{convección}, que necesitan materia; las
        superficies \textcolor{red}{plateadas} reflejan la \textcolor{red}{radiación} infrarroja. El tapón reduce las pérdidas por
        convección.
    \end{solu}
\end{exercise}

\begin{exercise}
    \(\star\star\) Un astronauta en el espacio debe mantener estable la temperatura de su cuerpo. Investiga cómo están hechos los
    trajes espaciales y qué mecanismos de transmisión del calor intervienen.
    \begin{solu}
        \textcolor{red}{Respuesta abierta}. En el vacío solo hay \textcolor{red}{radiación}: el traje tiene capas exteriores
        reflectantes y varias capas aislantes (aislamiento multicapa); para evacuar el calor del cuerpo usa una prenda interior con
        tubos por los que circula \textcolor{red}{agua de refrigeración}.
    \end{solu}
\end{exercise}

\subsection{Transferencia de energía mediante calor}

\begin{exercise}
    \(\star\) Define con tus palabras el calor específico. ¿En qué unidades se mide?
    \begin{solu}
        Es el \textcolor{red}{calor que necesita la unidad de masa de una sustancia para que su temperatura suba un grado}. Se mide
        en \textcolor{red}{\unit{\joule\per\kg\per\kelvin}} en el SI, o en \unit{\calorie\per\g\per\degreeCelsius}.
    \end{solu}
\end{exercise}

\begin{exercise}
    \(\star\) El calor específico del etanol es \qty{0.58}{\calorie\per\g\per\degreeCelsius}. Exprésalo en unidades del SI.
    \begin{solu}
        \(\num{0.58} \cdot \dfrac{\qty{4.18}{\joule}}{\qty{e-3}{\kg} \cdot \qty{1}{\kelvin}} = \textcolor{red}{\qty{2.42e3}{\joule\per\kg\per\kelvin}}\).
    \end{solu}
\end{exercise}

\begin{exercise}
    \(\star\star\) Una sustancia A tiene el doble de masa y el doble de calor específico que otra B. Si les suministramos el mismo
    calor, ¿cuál sufrirá un mayor aumento de temperatura?
    \begin{solu}
        \(\Delta T_A = \dfrac{Q}{2m \cdot 2c} = \dfrac{\Delta T_B}{4}\): \textcolor{red}{B se calienta cuatro veces más}.
    \end{solu}
\end{exercise}

\begin{exercise}
    \(\star\star\) Un bloque de cobre de \qty{1.2}{\kg} está a \qty{20}{\degreeCelsius} y le suministramos \qty{40}{\kilo\joule}. ¿Qué
    temperatura alcanza? \(c_{\text{Cu}} = \qty{0.093}{\calorie\per\g\per\degreeCelsius}\).
    \begin{solu}
        \(c = \qty{389}{\joule\per\kg\per\kelvin}\); \(\Delta T = \dfrac{40000}{\num{1.2} \cdot 389} = \qty{85.7}{\kelvin}\);
        \(T_f = \textcolor{red}{\qty{105.7}{\degreeCelsius}}\).
    \end{solu}
\end{exercise}

\begin{exercise}
    \(\star\star\) El recocido es una técnica de joyería que ablanda los metales para trabajarlos mejor. ¿Qué calor hay que dar a una
    pieza de plata de \qty{200}{\g} para calentarla desde \qty{25}{\degreeCelsius} hasta \qty{600}{\degreeCelsius}?
    \(c_{\text{Ag}} = \qty{0.056}{\calorie\per\g\per\degreeCelsius}\).
    \begin{solu}
        \(Q = 200 \cdot \num{0.056} \cdot 575 = \textcolor{red}{\qty{6440}{\calorie}} = \textcolor{red}{\qty{26.9}{\kilo\joule}}\).
    \end{solu}
\end{exercise}

\begin{exercise}
    \(\star\star\) Se mezclan \qty{4}{\kg} de agua a \qty{70}{\degreeCelsius} con \qty{12}{\kg} de agua a \qty{30}{\degreeCelsius}.
    Calcula la temperatura final y el calor que cede y absorbe cada muestra.
    \begin{solu}
        \(T = \dfrac{4 \cdot 70 + 12 \cdot 30}{16} = \textcolor{red}{\qty{40}{\degreeCelsius}}\); \(Q = 4 \cdot 4180 \cdot (40 - 70)
        = \textcolor{red}{\qty{-502}{\kilo\joule}}\) (cede la caliente) y \(\textcolor{red}{\qty{+502}{\kilo\joule}}\) (absorbe la fría).
    \end{solu}
\end{exercise}

\begin{exercise}
    \(\star\star\) Si se mezclan \qty{8}{\kg} de agua a \qty{50}{\degreeCelsius} con \qty{12}{\kg} de agua a \qty{15}{\degreeCelsius},
    ¿cuál será la temperatura final?
    \begin{solu}
        \(T = \dfrac{8 \cdot 50 + 12 \cdot 15}{20} = \textcolor{red}{\qty{29}{\degreeCelsius}}\).
    \end{solu}
\end{exercise}

\begin{exercise}
    \(\star\star\) ¿Qué calor hay que extraer de \qty{500}{\g} de agua a \qty{20}{\degreeCelsius} para convertirla en cubitos de
    hielo a \qty{-5}{\degreeCelsius}? Datos: \(L_f = \qty{80}{\calorie\per\g}\); \(c_{\text{hielo}} =
    \qty{0.50}{\calorie\per\g\per\degreeCelsius}\).
    \begin{solu}
        \(|Q| = 500\,(1 \cdot 20 + 80 + \num{0.5} \cdot 5) = \textcolor{red}{\qty{51250}{\calorie}} \approx \qty{214}{\kilo\joule}\).
        Con el criterio de signos, \(Q = \qty{-214}{\kilo\joule}\).
    \end{solu}
\end{exercise}

\begin{exercise}
    \(\star\star\) ¿Qué energía hace falta para fundir \qty{2}{\kg} de plomo que están a \qty{27}{\degreeCelsius}? Usa los datos de las
    tablas~\ref{tab:ud13-calor-especifico} y~\ref{tab:ud13-latentes}.
    \begin{solu}
        Calentar hasta \qty{327}{\degreeCelsius}: \(2 \cdot 130 \cdot 300 = \qty{78}{\kilo\joule}\); fundir: \(2 \cdot 25 =
        \qty{50}{\kilo\joule}\). Total: \(\textcolor{red}{\qty{128}{\kilo\joule}}\).
    \end{solu}
\end{exercise}

\begin{exercise}
    \(\star\star\) Un puente de hierro de \qty{500}{\m} tiene su longitud nominal a \qty{18}{\degreeCelsius}. Un día de verano, su
    longitud ha aumentado \qty{12}{\cm}. ¿Qué temperatura tiene? \(\lambda = \qty{1.2e-5}{\per\kelvin}\).
    \begin{solu}
        \(\Delta T = \dfrac{\Delta L}{L_0\,\lambda} = \dfrac{\num{0.12}}{500 \cdot \num{1.2e-5}} = \qty{20}{\kelvin}\);
        \(T = \textcolor{red}{\qty{38}{\degreeCelsius}}\).
    \end{solu}
\end{exercise}

\begin{exercise}
    \(\star\) ¿Cómo se pueden obtener los coeficientes de dilatación superficial y cúbica de un sólido a partir del lineal? Pon un
    ejemplo.
    \begin{solu}
        \(\beta \approx 2\lambda\) y \(\gamma \approx 3\lambda\). Para el acero, \(\lambda = \qty{1.2e-5}{\per\kelvin}\):
        \(\textcolor{red}{\beta \approx \qty{2.4e-5}{\per\kelvin}}\) y \(\textcolor{red}{\gamma \approx \qty{3.6e-5}{\per\kelvin}}\).
    \end{solu}
\end{exercise}

\subsection{Transferencia de energía mediante trabajo}

\begin{exercise}
    \(\star\) Define el trabajo termodinámico. ¿Es lo mismo un trabajo de expansión que uno de compresión?
    \begin{solu}
        Es la \textcolor{red}{energía intercambiada sin que intervenga una diferencia de temperatura}, en un gas al cambiar su
        volumen. \textcolor{red}{No}: en la expansión el gas hace trabajo sobre el entorno (\(W < 0\)); en la compresión el entorno lo
        hace sobre el gas (\(W > 0\)).
    \end{solu}
\end{exercise}

\begin{exercise}
    \(\star\star\) Un gas ocupa \qty{12}{\litre} y se expande hasta el doble de su volumen a una presión constante de \qty{3}{\atm}.
    Calcula el trabajo.
    \begin{solu}
        \(W = -3 \cdot 101325 \cdot \num{12e-3} = \textcolor{red}{\qty{-3.65}{\kilo\joule}}\) (\(\num{-36}~\mathrm{atm\cdot L}\)).
    \end{solu}
\end{exercise}

\begin{exercise}
    \(\star\star\) En el tiempo de compresión de un motor de gasolina, el volumen del cilindro pasa de \qty{450}{\cm\cubed} a
    \qty{40}{\cm\cubed}. Si suponemos una presión media constante de \qty{9}{\atm}, ¿qué trabajo se hace sobre el gas?
    \begin{solu}
        \(W = -9 \cdot 101325 \cdot (40 - 450) \cdot \num{e-6} = \textcolor{red}{\qty{+374}{\joule}}\).
    \end{solu}
\end{exercise}

\begin{exercise}
    \(\star\star\) ¿Qué trabajo realizan \qty{0.20}{\mol} de argón a \qty{30}{\degreeCelsius} al expandirse de forma isotérmica de
    \qty{2.0}{\litre} a \qty{4.5}{\litre}?
    \begin{solu}
        \(W = -\num{0.20} \cdot \num{8.31} \cdot \num{303} \cdot \ln\dfrac{\num{4.5}}{\num{2.0}} = \textcolor{red}{\qty{-409}{\joule}}\).
    \end{solu}
\end{exercise}

\begin{exercise}
    \(\star\star\) Un gas pasa del estado A (\qty{8}{\litre}, \qty{60}{\kilo\pascal}) al B (\qty{8}{\litre}, \qty{250}{\kilo\pascal})
    y después al C (\qty{20}{\litre}, \qty{250}{\kilo\pascal}). Representa el proceso en un diagrama \(p\)-\(V\) y calcula el trabajo
    total.
    \begin{solu}
        A \(\rightarrow\) B es isocórico (\(W = 0\)) y B \(\rightarrow\) C, isobárico:
        \(W = -\num{2.5e5} \cdot (20 - 8) \cdot \num{e-3} = \textcolor{red}{\qty{-3000}{\joule}}\).
    \end{solu}
\end{exercise}

\begin{exercise}
    \(\star\star\) Un gas ocupa \qty{4}{\litre} a \qty{2}{\atm} y recorre este ciclo: a) se calienta a presión constante hasta
    triplicar su volumen; b) se enfría a volumen constante hasta que la presión se reduce a la mitad; c) se comprime a presión
    constante hasta el volumen inicial; d) se calienta a volumen constante hasta la presión inicial. Represéntalo en un diagrama
    \(p\)-\(V\) y calcula el trabajo del ciclo.
    \begin{solu}
        \(W_a = -2 \cdot 8 = \num{-16}~\mathrm{atm\cdot L}\); \(W_b = W_d = 0\);
        \(W_c = -1 \cdot (-8) = \num{+8}~\mathrm{atm\cdot L}\).
        \(W = \num{-8}~\mathrm{atm\cdot L} = \textcolor{red}{\qty{-811}{\joule}}\).
    \end{solu}
\end{exercise}

\begin{exercise}
    \(\star\star\) ¿Qué expresiones son correctas para un proceso isobárico? a) \(W = -p\,\Delta V\); b) \(W = \Delta U\); c) \(W =
    -p\,(V_i - V_f)\); d) \(\Delta U = Q_p - p\,\Delta V\).
    \begin{solu}
        \textcolor{red}{a y d}. La b solo vale en un proceso adiabático, y la c tiene el signo cambiado.
    \end{solu}
\end{exercise}

\begin{exercise}
    \(\star\star\) Un gas ideal sufre tres etapas: A \(\rightarrow\) B, compresión a presión constante; B \(\rightarrow\) C,
    calentamiento a volumen constante; C \(\rightarrow\) A, expansión isotérmica que lo devuelve al estado inicial. Dibuja el ciclo en
    un diagrama \(p\)-\(V\) e indica el signo de \(W\), \(\Delta U\) y \(Q\) en cada etapa.
    \begin{solu}
        A \(\rightarrow\) B: \(W > 0\); la temperatura baja (menor \(V\) a igual \(p\)), \(\Delta U < 0\) y \(Q < 0\).
        B \(\rightarrow\) C: \(W = 0\), \(\Delta U > 0\), \(Q > 0\). C \(\rightarrow\) A: \(W < 0\), \(\Delta U = 0\), \(Q > 0\). El
        ciclo se recorre en sentido horario y el \textcolor{red}{trabajo neto es negativo}.
    \end{solu}
\end{exercise}

\subsection{Primer principio de la termodinámica}

\begin{exercise}
    \(\star\star\) Explica con tus palabras el primer principio de la termodinámica y el concepto de energía interna.
    \begin{solu}
        La \textcolor{red}{energía interna} es la suma de las energías de todas las partículas del sistema. El
        \textcolor{red}{primer principio} dice que solo cambia si el sistema intercambia calor o trabajo: \(\Delta U = Q + W\); es
        la conservación de la energía.
    \end{solu}
\end{exercise}

\begin{exercise}
    \(\star\star\) Haz un esquema de cómo se aplica el primer principio en los procesos adiabático, isocórico, isobárico, isotérmico
    y cíclico.
    \begin{solu}
        Adiabático: \(\textcolor{red}{\Delta U = W}\). Isocórico: \(\textcolor{red}{\Delta U = Q_V}\). Isobárico:
        \(\textcolor{red}{\Delta U = Q_p - p\,\Delta V}\). Isotérmico (gas ideal) y cíclico: \(\textcolor{red}{\Delta U = 0}\), \(Q = -W\).
    \end{solu}
\end{exercise}

\begin{exercise}
    \(\star\star\) Razona si son verdaderas o falsas: a) En un proceso isocórico, \(\Delta U = W\). b) La energía interna aumenta
    siempre que el sistema cede calor. c) En un proceso isotérmico de un gas ideal, \(\Delta U = Q_p - p\,\Delta V\). d) La energía
    interna disminuye siempre que el sistema recibe trabajo. e) En un proceso isobárico, \(\Delta U = 0\).
    \begin{solu}
        Todas \textcolor{red}{falsas}. a) \(\Delta U = Q_V\). b) Ceder calor tiende a disminuirla. c) \(\Delta U = 0\). d) Recibir
        trabajo tiende a aumentarla. e) En general \(\Delta U = Q_p - p\,\Delta V \neq 0\).
    \end{solu}
\end{exercise}

\begin{exercise}
    \(\star\star\) ¿Es cierto que en un proceso cíclico \(W = -Q\)? Razona tu respuesta.
    \begin{solu}
        \textcolor{red}{Sí}: el sistema vuelve al estado inicial y, como \(U\) es función de estado, \(\Delta U = 0 = Q + W\).
    \end{solu}
\end{exercise}

\begin{exercise}
    \(\star\star\) Cuando se pincha la rueda de una bicicleta, el aire que se escapa parece más frío que el del exterior. ¿Por qué?
    \begin{solu}
        Al salir, el aire \textcolor{red}{se expande muy deprisa}, sin tiempo para intercambiar calor (proceso adiabático):
        \(\Delta U = W < 0\), su energía interna disminuye y \textcolor{red}{se enfría}.
    \end{solu}
\end{exercise}

\begin{exercise}
    \(\star\star\) ¿Cuál es la variación de energía interna de un sistema que recibe \qty{400}{\calorie} de calor y sobre el que se
    realiza un trabajo de \qty{250}{\joule}?
    \begin{solu}
        \(\Delta U = 400 \cdot \num{4.18} + 250 = \textcolor{red}{\qty{1922}{\joule}}\).
    \end{solu}
\end{exercise}

\begin{exercise}
    \(\star\star\) La energía interna de un sistema aumenta \qty{320}{\joule} mientras cede \qty{450}{\joule} de calor. ¿Qué trabajo
    realiza el entorno sobre el sistema?
    \begin{solu}
        \(W = \Delta U - Q = 320 - (-450) = \textcolor{red}{\qty{+770}{\joule}}\).
    \end{solu}
\end{exercise}

\begin{exercise}
    \(\star\star\) ¿Aumentará la energía interna de un gas si se le somete a una compresión adiabática? Justifícalo.
    \begin{solu}
        \textcolor{red}{Sí}: \(Q = 0\) y \(W > 0\) (compresión), así que \(\Delta U = W > 0\); el gas se calienta. Es lo que ocurre
        en el tiempo de compresión de un motor diésel.
    \end{solu}
\end{exercise}

\begin{exercise}
    \(\star\star\) Clausius, Carnot y Joule hicieron aportaciones fundamentales a la termodinámica, pero ninguno era químico ni físico
    «de formación» (Carnot era ingeniero militar y Joule, cervecero). Investiga sus contribuciones y debatid en clase la importancia
    de otras disciplinas en el desarrollo de la física y de la química.
    \begin{solu}
        \textcolor{red}{Respuesta abierta}. Carnot: rendimiento máximo de las máquinas térmicas (1824). Joule: equivalente mecánico
        del calor. Clausius: primer y segundo principios y el concepto de entropía.
    \end{solu}
\end{exercise}

\subsection{Segundo principio de la termodinámica}

\begin{exercise}
    \(\star\star\) Explica con tus palabras el segundo principio de la termodinámica y el concepto de entropía.
    \begin{solu}
        La \textcolor{red}{entropía} mide el desorden molecular de un sistema. El \textcolor{red}{segundo principio} dice que en todo
        proceso espontáneo la entropía del universo aumenta; por eso el calor pasa del cuerpo caliente al frío y no se puede convertir
        todo el calor en trabajo.
    \end{solu}
\end{exercise}

\begin{exercise}
    \(\star\star\) ¿Sería correcto decir que el segundo principio es una restricción del primero? Justifica tu respuesta.
    \begin{solu}
        En cierto modo, \textcolor{red}{sí}: el primer principio dice que la energía se conserva; el segundo añade que
        \textcolor{red}{no todos los procesos que conservan la energía son posibles}, solo los que no disminuyen la entropía del
        universo (el café no se calienta solo enfriando el aire, aunque se conservara la energía).
    \end{solu}
\end{exercise}

\begin{exercise}
    \(\star\) Justifica por qué el grafito tiene mayor entropía que el diamante.
    \begin{solu}
        En el diamante cada carbono está unido rígidamente a otros cuatro en una red tridimensional; el grafito está formado por
        \textcolor{red}{capas} unidas por fuerzas débiles que pueden deslizar: su estructura es \textcolor{red}{menos ordenada}.
    \end{solu}
\end{exercise}

\begin{exercise}
    \(\star\star\) En un capítulo de una serie de dibujos animados, una niña construye una máquina que se mueve sin parar y sin
    consumir nada, y su padre exclama: «¡En esta casa obedecemos las leyes de la termodinámica!». ¿Tiene razón? ¿Qué principios
    violaría una máquina de movimiento perpetuo?
    \begin{solu}
        \textcolor{red}{Sí}. Una máquina que produjera trabajo sin consumir energía violaría el \textcolor{red}{primer principio}; una
        que convirtiera íntegramente en trabajo el calor de un solo foco (por ejemplo, el del aire) violaría el
        \textcolor{red}{segundo}.
    \end{solu}
\end{exercise}

\begin{exercise}
    \(\star\star\) Da una opinión razonada sobre esta afirmación: «Todo el trabajo mecánico puede convertirse en calor, pero no todo
    el calor puede convertirse en trabajo mecánico».
    \begin{solu}
        \textcolor{red}{Es correcta}: el trabajo se transforma íntegramente en calor (rozamiento), pero una máquina térmica siempre
        debe ceder parte del calor a un foco frío (\(\eta < 1\)), como exige el segundo principio.
    \end{solu}
\end{exercise}

\begin{exercise}
    \(\star\star\) Pronostica cómo cambia la entropía en estas reacciones: a) \ch{2 KClO3(s) -> 2 KCl(s) + 3 O2(g)}; b)
    \ch{2 CO(g) + O2(g) -> 2 CO2(g)}; c) \ch{Zn(s) + 2 HCl(aq) -> ZnCl2(aq) + H2(g)}; d) \ch{H2O(l) -> H2O(s)}; e) \ch{CH4(g) + 2
    O2(g) -> CO2(g) + 2 H2O(l)}.
    \begin{solu}
        a) \textcolor{red}{Aumenta} (se forma gas). b) \textcolor{red}{Disminuye} (de 3 a 2 moles de gas). c) \textcolor{red}{Aumenta}.
        d) \textcolor{red}{Disminuye}. e) \textcolor{red}{Disminuye} (de 3 a 1 mol de gas).
    \end{solu}
\end{exercise}

\begin{exercise}
    \(\star\star\) ¿Cuál es el rendimiento de una máquina térmica que realiza un trabajo de \qty{150}{\joule} por ciclo absorbiendo
    \qty{120}{\calorie}? ¿Qué calor cede en cada ciclo?
    \begin{solu}
        \(Q_1 = \qty{501.6}{\joule}\); \(\eta = \dfrac{150}{\num{501.6}} = \textcolor{red}{\qty{29.9}{\percent}}\);
        \(|Q_2| = \textcolor{red}{\qty{352}{\joule}}\).
    \end{solu}
\end{exercise}

\begin{exercise}
    \(\star\star\) Una máquina de Carnot trabaja entre \qty{400}{\kelvin} y \qty{300}{\kelvin} y absorbe \qty{1200}{\joule} del foco
    caliente en cada ciclo. Calcula su rendimiento, el trabajo que realiza y el calor que cede.
    \begin{solu}
        \(\eta = 1 - \dfrac{300}{400} = \textcolor{red}{\qty{25}{\percent}}\); \(|W| = \textcolor{red}{\qty{300}{\joule}}\);
        \(|Q_2| = \textcolor{red}{\qty{900}{\joule}}\).
    \end{solu}
\end{exercise}

%========================================================================================
\section{Ejercicios de ampliación}\label{sec:ud13-ampliacion}
%========================================================================================

\begin{exercise}
    \(\star\star\star\) Joule estableció que «en un sistema aislado térmicamente, el trabajo necesario para pasar de un estado A a
    un estado B es independiente del camino». Investiga sus experiencias y elabora un informe que incluya un ejemplo numérico: ¿qué
    aumento de temperatura sufriría el agua de una cascada de \qty{50}{\m} si toda su energía potencial se convirtiera en energía
    interna?
    \begin{solu}
        \textcolor{red}{Respuesta abierta}. Ejemplo: \(m\,g\,h = m\,c\,\Delta T \Rightarrow \Delta T = \dfrac{\num{9.8} \cdot 50}{4180} =
        \textcolor{red}{\qty{0.12}{\kelvin}}\). (Joule llegó a intentar medirlo en una cascada de los Alpes durante su luna de miel.)
    \end{solu}
\end{exercise}

\begin{exercise}
    \(\star\star\star\) ¿Qué energía se necesita para transformar \qty{60}{\g} de hielo a \qty{-15}{\degreeCelsius} en vapor de agua a
    \qty{100}{\degreeCelsius}, a presión constante? Datos: \(L_f = \qty{80}{\calorie\per\g}\); \(L_v = \qty{540}{\calorie\per\g}\);
    \(c_{\text{hielo}} = \qty{0.50}{\calorie\per\g\per\degreeCelsius}\); \(c_{\text{agua}} = \qty{1}{\calorie\per\g\per\degreeCelsius}\).
    \begin{solu}
        \(Q = 60\,(\num{0.5} \cdot 15 + 80 + 1 \cdot 100 + 540) = \textcolor{red}{\qty{43650}{\calorie}} = \qty{182}{\kilo\joule}\).
    \end{solu}
\end{exercise}

\begin{exercise}
    \(\star\star\star\) Un mol de hidrógeno a \qty{220}{\degreeCelsius} se enfría hasta \qty{100}{\degreeCelsius} a una presión
    constante de \qty{1}{\atm}. Considerándolo un gas ideal, calcula el calor y el trabajo intercambiados y la variación de energía
    interna. Dato: \(c_p(\ch{H2}) = \qty{14.3}{\joule\per\g\per\kelvin}\).
    \begin{solu}
        \(Q = \num{2.016} \cdot \num{14.3} \cdot (-120) = \textcolor{red}{\qty{-3.46}{\kilo\joule}}\).
        \(W = -p\,\Delta V = -n\,R\,\Delta T = -1 \cdot \num{8.31} \cdot (-120) = \textcolor{red}{\qty{+997}{\joule}}\).
        \(\Delta U = Q + W = \textcolor{red}{\qty{-2.46}{\kilo\joule}}\).
    \end{solu}
\end{exercise}

\begin{exercise}
    \(\star\star\star\) \qty{2.0}{\mol} de un gas ideal realizan un trabajo de \qty{4500}{\joule} sobre el entorno al expandirse de
    forma isotérmica hasta una presión de \qty{1.2}{\atm} y un volumen de \qty{40}{\litre}. Calcula: a) la temperatura del gas; b) su
    volumen inicial.
    \begin{solu}
        a) \(T = \dfrac{p\,V}{n\,R} = \dfrac{\num{1.2} \cdot 101325 \cdot \num{0.040}}{2 \cdot \num{8.31}} = \textcolor{red}{\qty{293}{\kelvin}}\).
        b) \(\ln\dfrac{V_f}{V_i} = \dfrac{4500}{n\,R\,T} = \num{0.925} \Rightarrow V_i = \dfrac{40}{e^{\num{0.925}}} =
        \textcolor{red}{\qty{15.9}{\litre}}\).
    \end{solu}
\end{exercise}

\begin{figure}[H]
    \centering
    \begin{tikzpicture}
        \begin{axis}[grafica, width=0.5\linewidth, height=5cm, xlabel={\(V\) (L)}, ylabel={\(p\) (atm)}, xmin=0, xmax=7.5,
            ymin=0, ymax=4.3, xtick={2,5,6}, ytick={1,2,3.5}]
            \addplot[very thick, ocre, mark=*, mark size=1.5pt] coordinates {(2,1) (2,2) (5,2) (5,3.5) (6,3.5)};
            \node[font=\small, anchor=west] at (axis cs:2,1) {A};
            \node[font=\small, anchor=south east] at (axis cs:2,2) {B};
            \node[font=\small, anchor=north west] at (axis cs:5,2) {C};
            \node[font=\small, anchor=south east] at (axis cs:5,3.5) {D};
            \node[font=\small, anchor=south west] at (axis cs:6,3.5) {E};
        \end{axis}
    \end{tikzpicture}
    \caption{Etapas de un gas en un diagrama \(p\)-\(V\).}
    \label{fig:ud13-amp-etapas}
\end{figure}

\begin{exercise}
    \(\star\star\star\) Un gas recorre las etapas de la figura~\ref{fig:ud13-amp-etapas}: A \(\rightarrow\) B, B \(\rightarrow\) C,
    C \(\rightarrow\) D y D \(\rightarrow\) E. Calcula el trabajo en cada una y el total.
    \begin{solu}
        \(W_{AB} = W_{CD} = \textcolor{red}{0}\) (isocóricos); \(W_{BC} = -2 \cdot 101325 \cdot \num{3e-3} = \textcolor{red}{\qty{-608}{\joule}}\);
        \(W_{DE} = -\num{3.5} \cdot 101325 \cdot \num{1e-3} = \textcolor{red}{\qty{-355}{\joule}}\). Total: \(\textcolor{red}{\qty{-963}{\joule}}\).
    \end{solu}
\end{exercise}

\begin{exercise}
    \(\star\star\star\) Un gas pasa del estado A (\qty{1.5}{\litre}, \qty{2}{\atm}) al C (\qty{4}{\litre}, \qty{6}{\atm}) por dos
    caminos: camino 1, A \(\rightarrow\) B (\qty{1.5}{\litre}, \qty{6}{\atm}) \(\rightarrow\) C; camino 2, A \(\rightarrow\) D
    (\qty{4}{\litre}, \qty{2}{\atm}) \(\rightarrow\) C. a) Calcula el trabajo en cada camino. b) Si por el camino 1 el gas absorbe
    \qty{2600}{\joule} de calor, ¿cuánto vale \(\Delta U\)? ¿Qué calor absorbe por el camino 2? c) ¿Es el trabajo una función de
    estado?
    \begin{solu}
        a) \(W_1 = -6 \cdot 101325 \cdot \num{2.5e-3} = \textcolor{red}{\qty{-1520}{\joule}}\); \(W_2 = -2 \cdot 101325 \cdot \num{2.5e-3} =
        \textcolor{red}{\qty{-507}{\joule}}\). b) \(\Delta U = 2600 - 1520 = \textcolor{red}{\qty{1080}{\joule}}\), igual por los dos
        caminos; \(Q_2 = 1080 + 507 = \textcolor{red}{\qty{1587}{\joule}}\). c) \textcolor{red}{No}: depende del camino, mientras que
        \(\Delta U\) no.
    \end{solu}
\end{exercise}

\begin{exercise}
    \(\star\star\star\) Una bala de plomo de \qty{10}{\g} que viaja a \qty{400}{\m\per\s} se incrusta en un bloque de madera. Si el
    \qty{60}{\percent} de su energía cinética se invierte en aumentar la energía interna de la bala, ¿cuánto sube su temperatura?
    Si estaba a \qty{20}{\degreeCelsius}, ¿llegará a fundirse? Usa los datos de las tablas de la unidad.
    \begin{solu}
        \(\Delta T = \dfrac{\num{0.6} \cdot \frac{1}{2}\,v^2}{c} = \dfrac{\num{0.6} \cdot 80000}{130} = \textcolor{red}{\qty{369}{\kelvin}}\).
        Llegaría a \qty{389}{\degreeCelsius} \(>\) \qty{327}{\degreeCelsius}: \textcolor{red}{empezaría a fundirse} (en parte, porque
        también hay que aportar el calor latente).
    \end{solu}
\end{exercise}

\begin{exercise}
    \(\star\star\star\) Un frigorífico extrae \qty{300}{\joule} de calor de su interior en cada ciclo y el compresor consume
    \qty{100}{\joule} de trabajo. a) ¿Qué calor cede a la cocina en cada ciclo? b) Su eficiencia es el cociente entre el calor extraído
    y el trabajo consumido; calcúlala. c) Si el interior está a \qty{4}{\degreeCelsius} y la cocina a \qty{25}{\degreeCelsius}, la
    eficiencia máxima es \(T_2/(T_1 - T_2)\). ¿Cuánto vale?
    \begin{solu}
        a) \(|Q_1| = 300 + 100 = \textcolor{red}{\qty{400}{\joule}}\). b) \(\varepsilon = \dfrac{300}{100} = \textcolor{red}{3}\). c)
        \(\varepsilon_{\text{máx}} = \dfrac{277}{21} = \textcolor{red}{\num{13.2}}\).
    \end{solu}
\end{exercise}

\begin{figure}[H]
    \centering
    \begin{tikzpicture}
        \begin{axis}[grafica, width=0.33\linewidth, height=4cm, xlabel={\(V\)}, ylabel={\(p\)}, xmin=0, xmax=6, ymin=0, ymax=5,
            xtick=\empty, ytick=\empty, title={\small A}]
            \addplot[fill=ocre!25, draw=none, domain=1:5] {4/x} \closedcycle;
            \addplot[very thick, ocre, domain=1:5] {4/x};
        \end{axis}
    \end{tikzpicture}\hfill
    \begin{tikzpicture}
        \begin{axis}[grafica, width=0.33\linewidth, height=4cm, xlabel={\(V\)}, ylabel={\(p\)}, xmin=0, xmax=6, ymin=0, ymax=5,
            xtick=\empty, ytick=\empty, title={\small B}]
            \addplot[fill=ocre!25, draw=none] coordinates {(1,0) (1,4) (5,0.8) (5,0)} \closedcycle;
            \addplot[very thick, ocre] coordinates {(1,4) (5,0.8)};
        \end{axis}
    \end{tikzpicture}\hfill
    \begin{tikzpicture}
        \begin{axis}[grafica, width=0.33\linewidth, height=4cm, xlabel={\(V\)}, ylabel={\(p\)}, xmin=0, xmax=6, ymin=0, ymax=5,
            xtick=\empty, ytick=\empty, title={\small C}]
            \addplot[fill=ocre!25, draw=none] coordinates {(1,0) (1,4) (5,4) (5,0)} \closedcycle;
            \addplot[very thick, ocre] coordinates {(1,0.8) (1,4) (5,4)};
        \end{axis}
    \end{tikzpicture}
    \caption{Tres procesos entre los mismos volúmenes.}
    \label{fig:ud13-amp-diagramas}
\end{figure}

\begin{exercise}
    \(\star\star\star\) Determina qué diagrama de la figura~\ref{fig:ud13-amp-diagramas} corresponde a un mayor trabajo realizado
    por el gas (en valor absoluto), sabiendo que en los tres el gas pasa del mismo volumen inicial al mismo volumen final.
    \begin{solu}
        El trabajo es el área bajo la curva: \(|W_C| > |W_B| > |W_A|\). En C el gas se expande a la presión más alta; el mayor
        trabajo corresponde al \textcolor{red}{diagrama C}. Esto muestra, de nuevo, que el trabajo depende del camino.
    \end{solu}
\end{exercise}

\begin{exercise}
    \(\star\star\star\) El ciclo de Carnot, del que se deduce el segundo principio, consta de dos procesos isotérmicos y dos
    adiabáticos. Por grupos, preparad una presentación que lo explique, junto con otros ciclos importantes como el de Otto (motores
    de gasolina) y el de Brayton (turbinas de gas y motores de avión), con sus diagramas \(p\)-\(V\).
    \begin{solu}
        \textcolor{red}{Respuesta abierta}. Carnot: expansión isotérmica, expansión adiabática, compresión isotérmica y compresión
        adiabática. Otto: dos adiabáticas y dos isocóricas. Brayton: dos adiabáticas y dos isobáricas.
    \end{solu}
\end{exercise}

%========================================================================================
\section{Autoevaluación}\label{sec:ud13-autoevaluacion}
%========================================================================================

\begin{exercise}
    Justifica si son verdaderas o falsas: a) Un cuerpo caliente contiene mucho calor. b) En un proceso adiabático no cambia la
    energía interna. c) El trabajo de expansión de un gas es negativo. d) La entropía de un sistema nunca puede disminuir.
    \begin{solu}
        a) \textcolor{red}{F}: tiene una temperatura alta y mucha energía interna; el calor es energía en tránsito. b)
        \textcolor{red}{F}: \(\Delta U = W\). c) \textcolor{red}{V}. d) \textcolor{red}{F}: puede disminuir; la que no disminuye es la
        del universo.
    \end{solu}
\end{exercise}

\begin{exercise}
    ¿Qué calor hay que dar a \qty{250}{\g} de agua para calentarla de \qty{20}{\degreeCelsius} a \qty{100}{\degreeCelsius}?
    \begin{solu}
        \(Q = \num{0.25} \cdot 4180 \cdot 80 = \textcolor{red}{\qty{83.6}{\kilo\joule}}\).
    \end{solu}
\end{exercise}

\begin{exercise}
    Se mezclan \qty{2}{\kg} de agua a \qty{60}{\degreeCelsius} con \qty{3}{\kg} de agua a \qty{10}{\degreeCelsius} en un recipiente
    aislado. ¿Cuál es la temperatura final?
    \begin{solu}
        \(T = \dfrac{2 \cdot 60 + 3 \cdot 10}{5} = \textcolor{red}{\qty{30}{\degreeCelsius}}\).
    \end{solu}
\end{exercise}

\begin{exercise}
    Un raíl de acero mide \qty{30}{\m} a \qty{0}{\degreeCelsius}. ¿Cuánto se alarga a \qty{40}{\degreeCelsius}? \(\lambda =
    \qty{1.2e-5}{\per\kelvin}\).
    \begin{solu}
        \(\Delta L = 30 \cdot \num{1.2e-5} \cdot 40 = \textcolor{red}{\qty{14.4}{\mm}}\).
    \end{solu}
\end{exercise}

\begin{exercise}
    Un gas se expande de \qty{2}{\litre} a \qty{5}{\litre} a una presión constante de \qty{100}{\kilo\pascal} y absorbe \qty{800}{\joule}
    de calor. Calcula el trabajo y la variación de energía interna.
    \begin{solu}
        \(W = -\num{1e5} \cdot \num{3e-3} = \textcolor{red}{\qty{-300}{\joule}}\); \(\Delta U = 800 - 300 = \textcolor{red}{\qty{500}{\joule}}\).
    \end{solu}
\end{exercise}

\begin{exercise}
    Un motor con un rendimiento del \qty{25}{\percent} absorbe \qty{4000}{\joule} de calor en cada ciclo. ¿Qué trabajo realiza y qué
    calor cede?
    \begin{solu}
        \(|W| = \num{0.25} \cdot 4000 = \textcolor{red}{\qty{1000}{\joule}}\); \(|Q_2| = \textcolor{red}{\qty{3000}{\joule}}\).
    \end{solu}
\end{exercise}

\begin{exercise}
    Indica si aumenta o disminuye la entropía: a) se evapora un charco; b) \ch{2 SO2(g) + O2(g) -> 2 SO3(g)}; c) se disuelve sal en
    agua.
    \begin{solu}
        a) \textcolor{red}{Aumenta}. b) \textcolor{red}{Disminuye} (de 3 a 2 moles de gas). c) \textcolor{red}{Aumenta}.
    \end{solu}
\end{exercise}
